% Propietario conservado: /Users/ruben/Documents/New project/output/EXTREMA_MEDIA_RAZON_RECIPROCIDAD_ES_EN_20260918/ARTICULO/ES/source/libro/04_bilateral.tex
\chapter{Conservation bilatérale et transduction arithmético-géométrique}
\label{x_reciprocidad:lib04:bilateral}

La constante holographique de couplage arithmético-géométrique désigne
le registre engendré \(K\), avec ses fenêtres ordonnées, son origine de phase,
son orientation et ses incidences. Sa production par le registre signé et
l’inversion de Hadamard est développée dans la \cref{x_reciprocidad:sec:k-registro}.
La lecture rationnelle \(\kappa\), le repère orbital \(O_K\) et les
opérateurs qui transportent ses réalisations sont des objets différents.
Cette distinction permet de composer leurs fonctions : génération du registre,
encodage arithmétique, transport de l’information et récupération des
lectures géométriques.

Les coordonnées régionales et \(\alpha\) sont reçues de la généalogie APP--TRIT--TPK, du même état enrichi et de sa structure discrète conjointe du continu. Les constructions suivantes agissent sur ces réalisations ultérieures. Leur contenu supplémentaire est une conservation bilatérale indépendante de la dimension, une extension réversible qui admet une mémoire entrante et une politique d'archivage avec récupération uniformément contrôlée.

\section{Loi bilatérale pour deux transports isométriques}
\label{x_reciprocidad:lib04:ley}

Soient \(H,H'\) des espaces de Hilbert et
\(\mathcal B,\mathcal C:H\to H'\) des isométries linéaires, de sorte que
\(\mathcal B^*\mathcal B=\mathcal C^*\mathcal C=I_H\).
La typographie calligraphique distingue ces transports de la base
arithmétique du raffinement. Pour \(a>0\), nous définissons
\begin{equation}
 \begin{aligned}
 p&=a^2+a+1=a(a+1)+1,\\
 N&=(2a+1)p=(a+1)^3+a^3,\\
 Q_-&=a\mathcal B-\mathcal C,\qquad
 Q_+=(a+1)\mathcal B+\mathcal C.
 \end{aligned}
 \label{x_reciprocidad:lib04:def-q}
\end{equation}

\begin{theorem}[Bilan de deux réalisations isométriques]
\label{x_reciprocidad:lib04:teo-balance}
Pour les applications \eqref{x_reciprocidad:lib04:def-q},
\begin{equation}
 (a+1)Q_-^*Q_-+aQ_+^*Q_+=NI_H.
 \label{x_reciprocidad:lib04:balance}
\end{equation}
L'identité vaut en toute dimension, y compris dans les espaces de Hilbert infinis, et admet des transports rectangulaires avec domaine et codomaine communs.
\end{theorem}

\begin{proof}
Écrivons \(X=\mathcal B^*\mathcal C+\mathcal C^*\mathcal B\).
L’isométrie des deux transports donne
\begin{align*}
 Q_-^*Q_-&=(a^2+1)I-aX,\\
 Q_+^*Q_+&=((a+1)^2+1)I+(a+1)X.
\end{align*}
Les pondérations \(a+1\) et \(a\) annulent exactement les termes
de \(X\). Le coefficient restant est
\[
 (a+1)(a^2+1)+a((a+1)^2+1)
 =2a^3+3a^2+3a+1=N.
\]
Les compositions sont celles d’opérateurs bornés et le calcul utilise
uniquement leurs produits et leurs adjoints. La preuve s’applique donc
également en dimension infinie.
\end{proof}

Pour \(a=8\), reçu de la partition nonadique précédente,
\begin{equation}
 9\|Q_-v\|^2+8\|Q_+v\|^2
 =(9^3+8^3)\|v\|^2=17\cdot73\,\|v\|^2.
 \label{x_reciprocidad:lib04:balance-canonico}
\end{equation}
Le facteur \(73=8\cdot9+1\) normalise cette identité générale. Dans le module ternaire, il se réalise en outre comme déterminant et facteur de similitude ; chacune de ces fonctions conserve ses propres hypothèses. Le bilan \eqref{x_reciprocidad:lib04:balance} utilise les isométries et admet aussi bien des transports d'ordre fini que des transports dépourvus de cette propriété.

\section{Factorisation, récupération et compatibilité des canaux}
\label{x_reciprocidad:lib04:factorizacion}

Définissons
\begin{equation}
 E_av=\frac1{\sqrt p}
 \begin{pmatrix}\sqrt{a(a+1)}\,\mathcal Bv\\\mathcal Cv\end{pmatrix},
 \qquad
 R_a=\frac1{\sqrt{2a+1}}
 \begin{pmatrix}\sqrt a&-\sqrt{a+1}\\
                 \sqrt{a+1}&\sqrt a\end{pmatrix}.
 \label{x_reciprocidad:lib04:factorizacion-er}
\end{equation}
La matrice \(R_a\) agit sur les deux canaux et est entendue comme tensorisée
avec \(I_{H'}\). On vérifie
\(E_a^*E_a=I_H\), \(R_a^*R_a=I_{H'\oplus H'}\), et
\begin{equation}
 W_av=R_aE_av
 =\frac1{\sqrt N}
 \begin{pmatrix}\sqrt{a+1}\,Q_-v\\\sqrt a\,Q_+v\end{pmatrix},
 \qquad W_a^*W_a=I_H.
 \label{x_reciprocidad:lib04:w}
\end{equation}
Dans le cas \(a=8\), la première répartition a pour poids \(72/73\) et
\(1/73\) ; la seconde opération est une rotation entre les canaux.
Cette factorisation exprime le rôle de \(72+1\) dans la norme
totale, en le distinguant de la restitution linéaire de
\eqref{x_reciprocidad:lib03:cadena-restauracion}.

\begin{proposition}[Récupération bilatérale]
\label{x_reciprocidad:lib04:prop-recuperacion}
Les canaux non normalisés \(x=Q_-v\), \(y=Q_+v\) déterminent
\begin{equation}
 \mathcal Bv=\frac{x+y}{2a+1},\qquad
 \mathcal Cv=\frac{ay-(a+1)x}{2a+1},\qquad
 v=\mathcal B^*\frac{x+y}{2a+1}.
 \label{x_reciprocidad:lib04:inversa-canales}
\end{equation}
À partir des canaux normalisés, on récupère
\(v=W_a^*W_av\), avec \(\|W_a^*\|=1\) lorsque \(H\neq\{0\}\).
\end{proposition}

\begin{proof}
Additionner les deux définitions de \(Q_\pm\) élimine \(\mathcal Cv\). La combinaison \(aQ_+-(a+1)Q_-\) élimine \(\mathcal Bv\). La première isométrie possède un inverse à gauche \(\mathcal B^*\), ce qui démontre les formules. La dernière affirmation découle de \eqref{x_reciprocidad:lib04:w} : l'adjoint d'une isométrie a une norme égale à un.
\end{proof}

L’image \(\operatorname{im}W_a\) est fermée et son projecteur
orthogonal est \(W_aW_a^*\). Un couple normalisé est compatible avec une
entrée unique exactement lorsqu’il appartient à cette image ; sa composante
orthogonale quantifie le défaut de compatibilité.

Si \(\mathcal B=I\) et si \(\mathcal C=U\) est unitaire, les deux
\(Q_\pm\) sont des polynômes du même opérateur. Pour les canaux non
normalisés, le critère se simplifie en
\begin{equation}
 Q_+x=Q_-y.
 \label{x_reciprocidad:lib04:compatibilidad-unitaria}
\end{equation}
En effet, l’égalité équivaut à
\(U(x+y)=ay-(a+1)x\). En définissant
\(v=(x+y)/(2a+1)\), on obtient par substitution
\(Q_-v=x\), \(Q_+v=y\). La nécessité résulte de la commutation
des deux polynômes. Cette forme simplifiée appartient au domaine
unitaire indiqué ; le critère d’image couvre les transports
généraux entre espaces distincts.

\section{Transformation des lecteurs et symétrie relative}
\label{x_reciprocidad:lib04:lectores}

\begin{theorem}[Lecteur commun dans les deux canaux]
\label{x_reciprocidad:lib04:teo-lector}
Pour tout opérateur borné autoadjoint \(F\) sur \(H'\),
\begin{equation}
 W_a^*\operatorname{diag}(F,F)W_a
 =\frac{a(a+1)\mathcal B^*F\mathcal B+
                  \mathcal C^*F\mathcal C}{a(a+1)+1}.
 \label{x_reciprocidad:lib04:lector-general}
\end{equation}
Si \(\mathcal B\) est unitaire et \(\mathcal C=\mathcal BR\), avec \(R\) unitaire, l'expression, pour \(G=\mathcal B^*F\mathcal B\), est
\begin{equation}
 \mathcal T_a(G)=\frac{a(a+1)G+R^*GR}{a(a+1)+1}.
 \label{x_reciprocidad:lib04:lector-relativo}
\end{equation}
On a \(\mathcal T_a(G)=G\) si et seulement si \(GR=RG\).
\end{theorem}

\begin{proof}
Le développement de \((a+1)Q_-^*FQ_-+aQ_+^*FQ_+\) annule les termes \(\mathcal B^*F\mathcal C+\mathcal C^*F\mathcal B\). Les coefficients des termes restants sont \(a(a+1)(2a+1)\) pour \(\mathcal B^*F\mathcal B\) et \(2a+1\) pour \(\mathcal C^*F\mathcal C\). Diviser par \(N\) démontre \eqref{x_reciprocidad:lib04:lector-general}. Substituer \(\mathcal C=\mathcal BR\) produit \eqref{x_reciprocidad:lib04:lector-relativo}. L'égalité avec \(G\) équivaut à \(R^*GR=G\), et l'unitarité de \(R\) la rend équivalente à \(GR=RG\).
\end{proof}

Le lecteur \(F=I\) récupère la norme. Pour un lecteur général,
l’identité détermine sa transformation, et la coïncidence
\(\mathcal B^*F\mathcal B=\mathcal C^*F\mathcal C=G\) garantit
sa conservation dans tous les états. Dans le cas canonique,
\begin{equation}
 \mathcal T_8(G)=\frac{72G+R^*GR}{73}.
 \label{x_reciprocidad:lib04:lector-72}
\end{equation}
Cette formule relie conservation et symétrie au moyen d’un commutateur
concret. L’observable qui transporte exactement un opérateur
\(A\) du domaine vers l’image complète est \(W_aAW_a^*\) ; son
identité ambiante est \(W_aW_a^*\). L’observable diagonal de
\eqref{x_reciprocidad:lib04:lector-general} est un autre choix, dont la différence est
calculée par le même théorème.

\section{Actualisation réversible avec mémoire entrante}
\label{x_reciprocidad:lib04:reentrada}

Soient maintenant \(\mathcal B=I\), \(\mathcal C=U\), avec \(U\)
unitaire dans un même espace, et écrivons
\begin{equation}
 A=\frac{\sqrt{a+1}}{\sqrt N}Q_-,\qquad
 D=\frac{\sqrt a}{\sqrt N}Q_+.
 \label{x_reciprocidad:lib04:ad}
\end{equation}
Les opérateurs \(A,D\) sont normaux et commutent entre eux et avec leurs adjoints, car ce sont des polynômes en \(U\) et \(U^*\).

\begin{theorem}[Extension unitaire]
\label{x_reciprocidad:lib04:teo-unitaria}
L’actualisation
\begin{equation}
 \mathscr R_a=
 \begin{pmatrix}A&-D^*\\D&A^*\end{pmatrix}:H\oplus H\to H\oplus H
 \label{x_reciprocidad:lib04:unitaria}
\end{equation}
est unitaire. Sa première colonne est \(W_a\), et une entrée
\((v,m)\) conserve \(\|v\|^2+\|m\|^2\).
\end{theorem}

\begin{proof}
Le bilan donne \(A^*A+D^*D=I\). La normalité fournit
\(AA^*+DD^*=I\). Les produits croisés de
\(\mathscr R_a^*\mathscr R_a\) et
\(\mathscr R_a\mathscr R_a^*\) s’annulent par les commutations
précédentes ; leurs deux diagonales sont l’identité. La première colonne
reproduit \eqref{x_reciprocidad:lib04:w}, et la conservation de la norme conjointe
est l’identité unitaire appliquée à \((v,m)\).
\end{proof}

Pour des données \(z=(z_1,z_2)\), l'inversion explicite est
\begin{equation}
 v_{\mathrm{rec}}=A^*z_1+D^*z_2,\qquad
 m_{\mathrm{rec}}=-Dz_1+Az_2.
 \label{x_reciprocidad:lib04:inversa-unitaria}
\end{equation}
L’hypothèse de mémoire entrante nulle se vérifie au moyen de
\(m_{\mathrm{rec}}=0\). Si \(e\) est une perturbation des données,
l’unitarité donne
\begin{equation}
 \|e\|^2=\|e_{\mathrm{rec}}\|^2+\|e_{\mathrm{mem}}\|^2.
 \label{x_reciprocidad:lib04:descomposicion-ruido}
\end{equation}
La seconde composante détecte la partie incompatible avec cette entrée. Une erreur appartenant à \(\operatorname{im}W_a\) a un résidu de mémoire nul et altère l'entrée récupérée. La correction du registre entier utilise en outre ses distances discrètes de séparation.


% BEGIN OWNER /Users/ruben/Documents/New project/output/EXTREMA_MEDIA_RAZON_RECIPROCIDAD_ES_EN_20260918/ARTICULO/ES/source/libro/memoria_reducida_20260916.tex
\section{Réduction causale de la mise à jour bilatérale}
\label{x_reciprocidad:lib04:reduccion-causal}

La mise à jour unitaire précédente réalise le transport de l’état et de sa
mémoire dans un même espace. L’observation de la première composante admet
une description équivalente qui conserve explicitement la contribution du
registre. Cette description s’obtient à partir du transport déjà construit, avec les
mêmes coefficients et la même orientation ; la mémoire éliminée de la liste
des variables réapparaît comme dépendance à l’égard des états précédents.

Fixons les opérateurs \(A,D\) de \eqref{x_reciprocidad:lib04:ad} et une réalisation \(H\oplus H\). Pour \(n\geq0\), soient
\begin{equation}
 v_{n+1}=Av_n-D^*m_n,\qquad
 m_{n+1}=Dv_n+A^*m_n.
 \label{x_reciprocidad:lib04:recurrencia-memoria}
\end{equation}
La mise à jour conjointe est \(\mathscr R_a\) et conserve \(\|v_n\|^2+\|m_n\|^2\).

\begin{proposition}[Élimination exacte du registre auxiliaire]
\label{x_reciprocidad:lib04:prop-memoria-reducida}
Pour tout \(n\geq0\), la somme vide étant interprétée comme zéro,
\begin{align}
 m_n&=(A^*)^n m_0+
       \sum_{k=0}^{n-1}(A^*)^{n-1-k}Dv_k,
       \label{x_reciprocidad:lib04:memoria-explicita}\\
 v_{n+1}&=Av_n-D^*(A^*)^n m_0
       +\sum_{k=0}^{n-1}\mathscr K_{n-1-k}v_k,
       \label{x_reciprocidad:lib04:evolucion-reducida}\\
 \mathscr K_j&=-D^*(A^*)^jD\quad(j\geq0).
       \label{x_reciprocidad:lib04:nucleo-retardo}
\end{align}
L’équation réduite, avec \eqref{x_reciprocidad:lib04:memoria-explicita}, est
équivalente à la récurrence conjointe pour les mêmes données \((v_0,m_0)\).
\end{proposition}
\begin{proof}
La première identité vaut pour \(n=0\). Si elle vaut pour \(n\), la seconde équation de \eqref{x_reciprocidad:lib04:recurrencia-memoria} fournit
\[
 m_{n+1}=(A^*)^{n+1}m_0+
 \sum_{k=0}^{n-1}(A^*)^{n-k}Dv_k+Dv_n,
\]
qui est l’identité dans \(n+1\). La substituer, à la profondeur \(n\), dans
la première équation démontre \eqref{x_reciprocidad:lib04:evolucion-reducida}.
Réciproquement, si \(v_n\) satisfait cette dernière et si l’on définit \(m_n\)
par \eqref{x_reciprocidad:lib04:memoria-explicita}, la même différence de sommes prouve
\(m_{n+1}=Dv_n+A^*m_n\) ; l’autre équation se récupère par substitution.
\end{proof}

Le noyau \(\mathscr K_j\) décrit un transfert vers le registre, \(j\) mises à jour à l'intérieur de celui-ci et une rentrée dans la composante lue. Le terme \(-D^*(A^*)^nm_0\) conserve la mémoire initiale. Sa suppression correspond au cas particulier \(m_0=0\), et non à une identité du transport général. Dans \eqref{x_reciprocidad:lib04:evolucion-reducida} interviennent uniquement l'état actuel et les états antérieurs. La causalité de cette réduction provient de l'ordre de la récurrence conjointe.

La normalisation du bilan fournit en outre une estimation explicite.
Comme \(Q_-=aI-U\), on a
\(\|A\|^2\leq r_a=(a+1)^3/((a+1)^3+a^3)<1\), tandis que
\(D^*D\leq I\). Par conséquent,
\begin{equation}
 \|\mathscr K_j\|\leq r_a^{j/2},\qquad
 \sum_{j\geq0}\|\mathscr K_j\|
 \leq\frac{1}{1-\sqrt{r_a}}.
 \label{x_reciprocidad:lib04:cota-nucleo-retardo}
\end{equation}
Cette borne contrôle le poids de la mémoire dans l’équation réduite ; la norme
de l’état conjoint reste conservée par la mise à jour unitaire.

La proposition est la réalisation temporelle de la réduction de Schur de
\cref{x_reciprocidad:mem:prop-schur}. Dans les séries formelles
\(V(z)=\sum_{n\geq0}v_nz^n\) et
\(M(z)=\sum_{n\geq0}m_nz^n\), les récurrences donnent
\(V=v_0+zAV-zD^*M\) et \(M=m_0+zDV+zA^*M\). En éliminant \(M\),
\begin{equation}
 V(z)=\bigl[I-zA+z^2D^*(I-zA^*)^{-1}D\bigr]^{-1}
       \bigl[v_0-zD^*(I-zA^*)^{-1}m_0\bigr].
 \label{x_reciprocidad:lib04:schur-bilateral}
\end{equation}
Les inverses formels existent car leurs termes constants sont des identités ;
les séries des états et la résolvante conjointe convergent également pour
\(|z|<1\). La source initiale et l’opérateur effectif se transforment ensemble.
Un lecteur qui agit également sur \(m_n\) conserve sa contribution au moyen de
\eqref{x_reciprocidad:lib04:memoria-explicita}, conformément à \eqref{x_reciprocidad:mem:schur-lector}.

La réentrée décrite ici actualise un registre entrant. La politique
d’archivage qui suit conserve en revanche chaque complément dans une coordonnée
distincte. Les deux opérations possèdent leur propre loi de conservation. Le noyau
\(\mathscr K_j\), les propagateurs de Green et une force dimensionnelle sont
des objets de types différents : le premier élimine une composante d’état,
les seconds résolvent des équations de source avec des conditions de support, et
la force ajoute une réalisation dynamique et ses unités.

% END OWNER


\section{Registre non stationnaire et convergence uniforme vers zéro de la composante terminale}
\label{x_reciprocidad:lib04:archivo}

À chaque étape, soient \(H_n,H_{n+1}\) des espaces de Hilbert et
\(\mathcal B_n,\mathcal C_n:H_n\to H_{n+1}\) des isométries. Le
paramètre \(a>0\) demeure fixe. Soient \(A_n,D_n\) les blocs
normalisés de \eqref{x_reciprocidad:lib04:w} ; on continue par le premier et on
enregistre le second :
\begin{equation}
 R_0=I,\qquad R_{n+1}=A_nR_n,\qquad m_n=D_nR_nv.
 \label{x_reciprocidad:lib04:politica-archivo}
\end{equation}
Cette politique possède un domaine distinct de celui de la réintroduction de
mémoire de \eqref{x_reciprocidad:lib04:unitaria} ; chacune conserve sa mise à jour.

\begin{lemma}[Transfert uniforme]
\label{x_reciprocidad:lib04:lem-tasa}
À chaque étape,
\begin{equation}
 \|A_n\|^2\leq r_a:={\frac{(a+1)^3}{(a+1)^3+a^3}}<1,
 \quad D_n^*D_n\geq(1-r_a)I,
 \quad\|R_n\|^2\leq r_a^n.
 \label{x_reciprocidad:lib04:tasa}
\end{equation}
\end{lemma}

\begin{proof}
L'inégalité triangulaire donne \(\|a\mathcal B_n-\mathcal C_n\|\leq a+1\). Sa normalisation fournit la première borne. La différence \(N-(a+1)^3=a^3>0\) prouve \(r_a<1\). Le bilan local implique \(D_n^*D_n=I-A_n^*A_n\), et multiplier les bornes des canaux poursuivis produit l'estimation de \(R_n\).
\end{proof}

\begin{theorem}[Mémoire complète à une profondeur arbitraire]
\label{x_reciprocidad:lib04:teo-archivo}
L’application finie
\begin{equation}
 Z_Nv=(R_Nv,D_0v,D_1R_1v,\ldots,D_{N-1}R_{N-1}v)
 \label{x_reciprocidad:lib04:zn}
\end{equation}
est isométrique. L’archive infinie
\begin{equation}
 Z_\infty v=(D_nR_nv)_{n\geq0}
 \label{x_reciprocidad:lib04:zinfty}
\end{equation}
appartient à la somme hilbertienne des espaces de sortie et satisfait
\begin{equation}
 \|v\|^2=\sum_{n\geq0}\|m_n\|^2,\qquad
 v=\sum_{n\geq0}R_n^*D_n^*m_n,
 \qquad Z_\infty^*Z_\infty=I.
 \label{x_reciprocidad:lib04:recuperacion-infinita}
\end{equation}
La série d’inversion converge en norme et l’adjoint est de norme un.
\end{theorem}

\begin{proof}
L’identité \(A_n^*A_n+D_n^*D_n=I\), conjuguée par \(R_n\),
produit
\[
 R_n^*R_n-R_{n+1}^*R_{n+1}=R_n^*D_n^*D_nR_n.
\]
La sommation de zéro à \(N-1\) donne
\begin{equation}
 I=R_N^*R_N+\sum_{n<N}R_n^*D_n^*D_nR_n.
 \label{x_reciprocidad:lib04:telescopia}
\end{equation}
C'est l'identité de Gram de \(Z_N\). D'après \eqref{x_reciprocidad:lib04:tasa}, le premier terme tend vers zéro en norme d'opérateur. L'identité limite prouve la convergence quadratique de l'archive et son isométrie. L'adjoint de cette isométrie est borné ; les troncatures de toute archive de carré sommable convergent dans la somme hilbertienne, et leur image par l'adjoint prouve la convergence de la série de \eqref{x_reciprocidad:lib04:recuperacion-infinita}. L'appliquer aux données exactes récupère \(v\).
\end{proof}

Avec seulement les \(N\) premiers compléments, l’adjoint tronqué a pour
résidu exact
\begin{equation}
 v-\sum_{n<N}R_n^*D_n^*m_n=R_N^*R_Nv,
 \qquad
 \left\|v-\sum_{n<N}R_n^*D_n^*m_n\right\|
 \leq r_a^N\|v\|.
 \label{x_reciprocidad:lib04:sesgo-truncado}
\end{equation}
Si l’erreur conjointe des données est au plus \(\varepsilon\), la
borne totale est \(\varepsilon+r_a^N\|v\|\). Avec \(Z_N\), y compris
son terminal, ou avec \(Z_\infty\), l’adjoint retrouve avec une erreur au
plus égale à \(\varepsilon\).

La politique précédente épuise le terminal par la borne uniforme prouvée.
Les évolutions précédentes qui possèdent un terminal positif
\(A_\infty\neq0\) conservent dans leur archive la composante
\(A_\infty^{1/2}v\). La profondeur arbitraire et l’élimination du
terminal sont donc des conclusions différentes et maintiennent leurs
hypothèses respectives.

\begin{proposition}[Inverse des mémoires finies sans terminal]
\label{x_reciprocidad:lib04:prop-inversa-finita}
Pour \(N\geq1\), soit \(M_Nv=(D_nR_nv)_{n<N}\). Alors
\begin{equation}
 M_N^*M_N=I-R_N^*R_N\geq(1-r_a^N)I,
 \label{x_reciprocidad:lib04:gram-finito}
\end{equation}
et l’inverse des moindres carrés
\begin{equation}
 L_N=(I-R_N^*R_N)^{-1}M_N^*,\qquad
 L_NM_N=I,\qquad
 \|L_N\|\leq(1-r_a^N)^{-1/2}
 \label{x_reciprocidad:lib04:inversa-finita}
\end{equation}
retrouve exactement l’entrée avec des données exactes.
\end{proposition}

\begin{proof}
L’identité de Gram est \eqref{x_reciprocidad:lib04:telescopia} ; la borne provient de
\eqref{x_reciprocidad:lib04:tasa}. L’opérateur positif est inversible par sa borne
inférieure stricte. La composition démontre \(L_NM_N=I\), et
\(L_NL_N^*=(M_N^*M_N)^{-1}\) prouve la norme déclarée.
\end{proof}

Cet inverse fini élimine le biais de
\eqref{x_reciprocidad:lib04:sesgo-truncado} ; son amplification du bruit est
spécifiée par \eqref{x_reciprocidad:lib04:inversa-finita}. L’adjoint tronqué,
l’inverse fini et l’adjoint de l’archive complète sont ainsi trois
opérateurs aux comportements distincts.

Pour \(a=8\), la borne générale est \(r_a=729/1241\). Dans le cas \(\mathcal B=I\), \(\mathcal C=S^4\), avec \(S\) le cycle de douze positions, le projecteur fixe \(P_0=(I+S^4+S^8)/3\) donne
\begin{equation}
 A^*A=\frac{a+1}{(2a+1)p}
       \bigl((a-1)^2P_0+p(I-P_0)\bigr).
 \label{x_reciprocidad:lib04:gram-tasa-ternaria}
\end{equation}
Comme \(p-(a-1)^2=3a>0\), sa norme est
\((a+1)/(2a+1)\). Dans \(a=8\), on obtient \(9/17\) ; la valeur propre
du secteur fixe est \(441/1241\). Ces taux appartiennent aux
transports et à la politique déclarés.

Si les deux canaux sont raffinés dans un arbre, chaque niveau complet est également isométrique par composition. Additionner les normes de toutes les couches complètes compterait de façon répétée la même norme initiale. La conservation est formulée au moyen des applications de raffinement entre niveaux ou au moyen de l'archive des compléments de \eqref{x_reciprocidad:lib04:zinfty}.

\section{Formes géométriques et conservation de leur transport}
\label{x_reciprocidad:lib04:formas}

Soient \(b,b'\) des formes non dégénérées, symétriques ou alternées, sur
les espaces correspondants, et supposons
\(\mathcal B^*b'\mathcal B=\mathcal C^*b'\mathcal C=b\).
Le même développement bilinéaire démontre
\begin{equation}
 (a+1)Q_-^*b'Q_-+aQ_+^*b'Q_+=Nb.
 \label{x_reciprocidad:lib04:balance-formas}
\end{equation}
En effet, les termes mixtes
\(\mathcal B^*b'\mathcal C+\mathcal C^*b'\mathcal B\) s’annulent
avec les mêmes pondérations ; les termes purs ont pour somme \(Nb\).
La forme transportée est ainsi conservée sans exiger qu’elle soit positive.
Les bornes contractives et de bruit de la sous-section précédente utilisent,
en revanche, des normes de Hilbert positives. Une forme lorentzienne ou
alternée conserve \eqref{x_reciprocidad:lib04:balance-formas}, avec ses propres
interprétations géométriques, et ne remplace pas la positivité dans ces
inégalités.

\section{Corrélation des canaux et problème inverse de lecture}
\label{x_reciprocidad:lib04:correlacion}

Pour \(v\neq0\), définissons
\begin{equation}
 \chi(v)=\frac{\operatorname{Re}\langle\mathcal Bv,\mathcal Cv\rangle}
                      {\|v\|^2},\qquad -1\leq\chi(v)\leq1.
 \label{x_reciprocidad:lib04:chi}
\end{equation}
La borne est Cauchy--Schwarz appliquée aux deux images isométriques.
Les lectures normalisées sont
\begin{equation}
 s_-=a^2+1-2a\chi,\qquad
 s_+=(a+1)^2+1+2(a+1)\chi.
 \label{x_reciprocidad:lib04:s-pm}
\end{equation}
Pour \(a=8\),
\begin{equation}
 s_-=65-16\chi,\qquad s_+=82+18\chi,\qquad
 9s_-+8s_+=1241.
 \label{x_reciprocidad:lib04:s-canonico}
\end{equation}
Les deux normes individuelles contiennent une corrélation réelle commune. Le bilan est une relation exacte entre elles, et non deux déterminations arithmétiques indépendantes d'une constante.

Utiliser la sortie préalable \(\alpha_A\) dans l’essai inverse
\(s_-=10^4\alpha_A\) détermine la lecture complémentaire
\begin{equation}
 s_+=73+\frac98(73-10^4\alpha_A)
       \simeq73{,}029783595557.
 \label{x_reciprocidad:lib04:ensayo-alpha}
\end{equation}
L’égalité s’obtient par substitution dans
\eqref{x_reciprocidad:lib04:s-canonico}. La moyenne pondérée
\((9s_-+8s_+)/17\) reste exactement égale à \(73\).
La comparaison conserve la généalogie préalable de \(\alpha_A\) : c’est
une condition de l’essai inverse, et non un autre producteur indépendant de
cette sortie. Un changement unitaire peut modifier \(\chi\) et les deux
lectures tout en maintenant le bilan, même s’il change une relation particulière
d’ordre ternaire.

\section{Repère cyclique et transport des réponses}
\label{x_reciprocidad:lib04:marco}

Sur l'espace porteur du registre engendré, soient \(H=\mathbb R^{12}\), \((Sx)_i=x_{i+1}\), avec indices cycliques, et
\begin{equation}
 O_K=(K,SK,\ldots,S^{11}K),\qquad
 589609I\leq O_K^*O_K\leq39225169I.
 \label{x_reciprocidad:lib04:gram-k}
\end{equation}
Les bornes sont celles du repère cyclique du registre canonique,
démontrées dans la \cref{x_reciprocidad:lib01:marco-ciclico}.
Soit \(Z\) l’une des isométries complètes
\(Z_N,Z_\infty\), ou une composition avec les archives précédentes
qui conserve leurs terminaux et leurs compléments.

\begin{theorem}[Conservation du conditionnement du repère]
\label{x_reciprocidad:lib04:teo-marco}
Si \(x=O_K\beta\), la donnée \(z=Zx\) détermine
\begin{equation}
 \beta=O_K^{-1}Z^*z,\qquad
 \|\delta\beta\|\leq\frac{\|\delta z\|}{\sqrt{589609}}.
 \label{x_reciprocidad:lib04:inversa-marco}
\end{equation}
Les douze colonnes de \(ZO_K\) forment une base de son image, même lorsque l'archive possède une infinité de coordonnées.
\end{theorem}

\begin{proof}
L'égalité \(Z^*Z=I\) donne \((ZO_K)^*(ZO_K)=O_K^*O_K\). Les bornes de Gram sont conservées et \(O_K\) est inversible grâce à la borne inférieure stricte. Composer leurs inverses démontre la récupération, et la norme \(\|O_K^{-1}\|\leq589609^{-1/2}\), avec \(\|Z^*\|=1\), donne la stabilité. Le nombre de colonnes et leur indépendance déterminent la dimension de l'image, sans l'identifier à tout l'espace de stockage.
\end{proof}

Pour un opérateur \(H_0\) qui commute avec \(S\), la réponse
vectorielle complète \(y=H_0K\) détermine
\begin{equation}
 H_0=O_yO_K^{-1}.
 \label{x_reciprocidad:lib04:identificacion-filtro}
\end{equation}
En effet, la commutation implique
\(H_0O_K=(y,Sy,\ldots,S^{11}y)=O_y\). Si l’on reçoit uniquement
\(z=Zy+e\), on définit \(\widehat y=Z^*z\) et
\(\widehat H_0=O_{\widehat y}O_K^{-1}\). Alors
\begin{equation}
 \|\widehat H_0-H_0\|
 \leq\sqrt{\frac{12}{589609}}\,\|e\|.
 \label{x_reciprocidad:lib04:error-filtro}
\end{equation}
La matrice orbitale de \(\delta y\) a douze colonnes de norme
\(\|\delta y\|\), de sorte que sa norme opératorielle ne dépasse pas
\(\sqrt{12}\|\delta y\|\). Cela et
\eqref{x_reciprocidad:lib04:gram-k} prouvent \eqref{x_reciprocidad:lib04:error-filtro}.
Pour un opérateur général, le protocole requiert les douze réponses
\(H_0S^jK\) ; une réponse dans le cas cyclique signifie ses douze
composantes, non une mesure scalaire.

Dans l’image transportée, le décalage est
\(S_Z=ZSZ^*\) et vérifie \(S_Z^jZ=ZS^j\). Cette égalité
détermine l’avancement de phase dans l’archive ; une translation de cases
de l’espace ambiant ne représente cet avancement que lorsqu’elle coïncide avec
l’opérateur transporté.

\section{Lecture positionnelle et précision scalaire}
\label{x_reciprocidad:lib04:kappa}

Fixons \(B_0=1000\), \(d=B_0^{12}-1\), et la ligne
\begin{equation}
 \ell=\frac{(B_0^{11},B_0^{10},\ldots,1)}{d}.
 \label{x_reciprocidad:lib04:fila-posicional}
\end{equation}
Son extension linéaire agit sur \(H\), et dans le domaine des mots \(\kappa(K)=\ell K\). Le représentant du lecteur transporté est \(Z\ell^*\), de sorte que
\begin{equation}
 \kappa(K)=\langle Z\ell^*,ZK\rangle,\qquad
 |\delta\kappa|\leq\|\ell\|\,\|e\|.
 \label{x_reciprocidad:lib04:lectura-kappa}
\end{equation}
La somme géométrique fournit
\begin{equation}
 \|\ell\|^2
 =\frac{B_0^{12}+1}{(B_0^2-1)(B_0^{12}-1)}.
 \label{x_reciprocidad:lib04:norma-fila}
\end{equation}
En effet, le numérateur de la somme des carrés est
\((B_0^{24}-1)/(B_0^2-1)\), qui, divisé par
\((B_0^{12}-1)^2\), donne \eqref{x_reciprocidad:lib04:norma-fila}.

Avec seulement \(N\) compléments et l’adjoint tronqué, il apparaît en outre
une erreur au plus égale à \(\|\ell\|r_a^N\|K\|\). L’inverse fini
\eqref{x_reciprocidad:lib04:inversa-finita} élimine ce biais et conserve sa borne
explicite d’amplification du bruit.

Dans le domaine \(\{0,\ldots,999\}^{12}\), la lecture exacte
retrouve \(N_K=d\kappa\) et ses douze blocs par divisions
euclidiennes successives. Si
\begin{equation}
 |\widetilde\kappa-\kappa|<\frac1{2d},
 \label{x_reciprocidad:lib04:radio-escalar}
\end{equation}
l'arrondi de \(d\widetilde\kappa\) à l'entier le plus proche récupère \(N_K\) : sa distance à l'entier véritable est inférieure à une demi-unité. La base, la longueur, l'origine et l'orientation accompagnent cette représentation. Le seuil scalaire de \eqref{x_reciprocidad:lib04:radio-escalar} se distingue des bornes vectorielles des registres de mémoire.

\section{Incidence, charge et restitution exacte du registre}
\label{x_reciprocidad:lib04:incidencia}

Soient \(\mathbf1=(1,\ldots,1)^T\) et
\(\Pi=I-\mathbf1\mathbf1^*/12\). L’isométrie incidencielle de
la \cref{x_reciprocidad:nuc05:isometria} est
\begin{equation}
 Fv=(\mu(v),\mathcal I\Pi v/6),\qquad
 \mu(v)=\frac{\mathbf1^*v}{12},\qquad
 \|(\mu,y)\|_Y^2=12\mu^2+\|y\|^2.
 \label{x_reciprocidad:lib04:f-incidencia}
\end{equation}
Son codomaine est son image déclarée. Le gramien
\(\mathcal I^*\mathcal I=36I+30\mathbf1\mathbf1^*\) donne
\(\|\mathcal I\Pi v\|^2=36\|\Pi v\|^2\), et
\(12\mu(v)^2+\|\Pi v\|^2=\|v\|^2\), d’où \(F^*F=I\).
Ainsi, à partir de \(z=Zv+e\),
\begin{equation}
 \widehat{Fv}=FZ^*z,\qquad \|FZ^*e\|_Y\leq\|e\|.
 \label{x_reciprocidad:lib04:error-incidencia}
\end{equation}
L’incidence non normalisée \(\mathcal I\Pi v\) a une erreur au
plus égale à \(6\|e\|\) ; le registre signé \(H_{12}v\), dont la
normalisation vérifie \(H_{12}^*H_{12}=4I\), a une erreur au plus égale à
\(2\|e\|\).

La direction \(u=P_3\Pi K\) et le projecteur
\(P_{10}=\Pi-uu^*/\|u\|^2\), avec \(u\neq0\), se récupèrent
avec leurs lecteurs et leur carte préalables. La sélection absolue
\(\{j:K_j\geq729\}\) a une discontinuité en \(K_4=729\).
La restitution entière s’effectue avant d’appliquer ce seuil.

\begin{proposition}[Décodage après l'archive complète]
\label{x_reciprocidad:lib04:prop-redondeo}
Soit \(K\) entier et \(z=ZK+e\), avec \(Z^*Z=I\). Si \(\|e\|<1/2\), arrondir chaque coordonnée de \(Z^*z\) récupère \(K\) exactement. Par composition, on restitue \(\kappa\), les signatures et les lecteurs géométriques définis sur ce registre.
\end{proposition}

\begin{proof}
La contractivité de \(Z^*\) donne
\(\|Z^*z-K\|\leq\|e\|<1/2\). Chaque erreur coordonnée reste
bornée par la norme conjointe, de sorte que l’unique coordonnée entière
à une distance inférieure à une demi-unité est l’originale. Une fois
\(K\) retrouvé, l’évaluation de l’un de ses lecteurs exacts
restitue sa valeur.
\end{proof}

En particulier, on récupère le support \(\{4,6,9,10\}\), y compris
sa composante de frontière. Les autres marques qui interviennent dans un
drapeau exceptionnel sont conservées par leurs canaux propres : la
restitution de \(K\) corrige exactement l’information qui y trouve
ses arguments suffisants.

\section{Conservation du conditionnement de deux mémoires}
\label{x_reciprocidad:lib04:dos-memorias}

Maintenons les lecteurs antécédents
\begin{equation}
 T_j=\frac{8I+S^j}{9},\qquad
 D_j=\frac{\sqrt8}{9}(S^j-I),\qquad
 Mv=(D_3v,D_4T_3v).
 \label{x_reciprocidad:lib04:m-antecedente}
\end{equation}
Les \(D_j\) de cette équation sont les compléments précédents ; le second bloc bilatéral \(D\) de \eqref{x_reciprocidad:lib04:ad} conserve sa propre définition. On a \(\ker M=\mathbb R\mathbf1\) et un inverse \(L_M\) de norme \(9/4\) sur le secteur centré, comme démontré dans la \cref{x_reciprocidad:nuclear:11}.

Soit \(E=W_a\oplus W_a\) dans l’espace des deux mémoires.
Son isométrie implique
\begin{equation}
 (EM)^*(EM)=M^*M,\qquad
 \ker(EM)=\mathbb R\mathbf1.
 \label{x_reciprocidad:lib04:gram-m-preservado}
\end{equation}
Pour \(z=EMv+e\) et la charge exacte
\(q=\mathbf1^*v\), la reconstruction est
\begin{equation}
 \widehat v=\frac q{12}\mathbf1+L_ME^*z,\qquad
 \|\widehat v-v\|\leq\frac94\|e\|.
 \label{x_reciprocidad:lib04:inversa-m}
\end{equation}
La preuve consiste à appliquer \(E^*E=I\), puis à utiliser
\(L_MM=\Pi\) et \(\|E^*\|=1\). Si la charge a une erreur
\(\delta q\), les secteurs uniforme et centré sont orthogonaux et
\begin{equation}
 \|\widehat v-v\|^2
 \leq\frac{|\delta q|^2}{12}+\frac{81}{16}\|e\|^2.
 \label{x_reciprocidad:lib04:error-carga}
\end{equation}

Les distances de tous les mots de l’image entière sont conservées
exactement par \(E\). Les rayons de restitution de
la \cref{x_reciprocidad:lib02:radios-memoria} demeurent donc
\begin{equation}
 r_0=\frac{2\sqrt{146}}{81}\quad\hbox{sans charge},\qquad
 r_q=\frac{4\sqrt{73}}{81}\quad\hbox{avec charge exacte}.
 \label{x_reciprocidad:lib04:radios}
\end{equation}
À l'intérieur du premier rayon, on récupère \(\Pi K\), et avec lui la direction et son projecteur lorsqu'ils sont définis. À l'intérieur du second, en conservant \(q\), on récupère \(K\) et le sélecteur absolu de seuil. Composer ensuite une autre isométrie complète, finie ou infinie, maintient les mêmes garanties : on applique son adjoint avant le décodeur. Le rayon est toujours exprimé dans la norme conjointe des données avec la normalisation fixée.

Cette composition redistribue l’information de \(M\) sans détériorer
son conditionnement et sans créer le canal uniforme que \(M\) élimine.
La conservation du Gram distingue cette propriété d’une amélioration
numérique obtenue en changeant le lecteur ou l’échelle du bruit.

\section{Norme orbitale et normalisation polaire du repère}
\label{x_reciprocidad:lib04:normalizacion}

L’inversion précédente satisfait
\(K=H_{12}U_{\mathrm{sgn}}/4\),
\(H_{12}^*H_{12}=4I\). Par conséquent
\begin{equation}
 \|K\|^2=\frac14\|U_{\mathrm{sgn}}\|^2=4251257,
 \qquad \|U_{\mathrm{sgn}}\|^2=17005028.
 \label{x_reciprocidad:lib04:norma-k}
\end{equation}
Chaque translation \(S^j\) est orthogonale et conserve
\(\nu=\sqrt{4251257}\). D’où
\begin{equation}
 \operatorname{tr}(O_K^*O_K)=12\nu^2=51015084.
 \label{x_reciprocidad:lib04:traza-marco}
\end{equation}
La division par \(\nu\) normalise chaque colonne orbitale à la norme
un et conserve son ordre. Elle est uniforme dans l’orbite et sous ses
transports isométriques. Les registres distincts maintiennent leurs normes
propres, et les colonnes normalisées continuent d’avoir un gramien
non scalaire : ses valeurs propres sont les précédentes divisées par
\(4251257\).

La normalisation opératoire utilise le Gram complet :
\begin{equation}
 G_K=O_K^*O_K,\qquad
 F_K=O_KG_K^{-1/2},\qquad
 F_K^*F_K=I.
 \label{x_reciprocidad:lib04:polar}
\end{equation}
La racine inverse existe par \eqref{x_reciprocidad:lib04:gram-k}. L’égalité se
vérifie en multipliant par \(G_K^{-1/2}\) des deux côtés du Gram.
Ainsi, \(ZF_K\) est isométrique à toute profondeur.
La récupération intégrale du repère est
\begin{equation}
 O_K=F_KG_K^{1/2}.
 \label{x_reciprocidad:lib04:polar-inversa}
\end{equation}
La paire \((F_K,G_K)\) conserve l'information complète. Le facteur unitaire conserve un repère orthonormal orienté dans les cartes fixées, tandis que les amplitudes résident également dans \(G_K\). Multiplier \(K\) par un scalaire positif maintient \(F_K\) et change \(G_K\) ; \(K\) et \(-K\) partagent une matrice de Gram et possèdent des facteurs polaires opposés, avec une lecture linéaire distincte \(\kappa\). Chaque exemple montre pourquoi un facteur isolé ne remplace pas la paire.

La translation de l’origine conserve \(\nu\), mais dans la carte
positionnelle fixe
\begin{equation}
 \kappa(SK)=1000\kappa(K)-K_1.
 \label{x_reciprocidad:lib04:cambio-origen}
\end{equation}
Cela s’obtient en décalant le numérateur positionnel et en utilisant
\(d=1000^{12}-1\). Pour conserver la même lecture sous un changement
de carte, on transporte aussi la ligne \(\ell\) par l’inverse
de ce changement. Une carte unitaire générale peut quitter le réseau des
blocs entiers ; le décodage par divisions euclidiennes se réalise
dans la carte de blocs récupérée.

\section{Exemple exact de transduction et de correction d’erreurs}
\label{x_reciprocidad:lib04:ejemplo}

Prenons \(a=8\), \(U=S^4\) et le registre canonique
\begin{equation}
 \begin{split}
 K=(&234,543,140,729,659,824,\\
    &621,58,914,794,146,601).
 \end{split}
 \label{x_reciprocidad:lib04:k-ejemplo}
\end{equation}
En notant \((v)_{1:6}\) et \((v)_{7:12}\) ses deux blocs de
six coordonnées, les lectures non normalisées sont
\begin{align}
 (Q_-K)_{1:6}&=(1213,3520,499,5774,4358,5798),\nonumber\\
 (Q_-K)_{7:12}&=(4822,-137,7078,5809,1028,4079),\nonumber\\
 (Q_+K)_{1:6}&=(2765,5711,1881,6619,6845,8210),\nonumber\\
 (Q_+K)_{7:12}&=(5735,1123,8460,7689,1454,6138).
 \label{x_reciprocidad:lib04:canales-ejemplo}
\end{align}
Leur somme, divisée par \(17\), retrouve tous les blocs de \(K\).
L’identité matricielle
\(9Q_-^*Q_-+8Q_+^*Q_+=1241I\) précède son évaluation sur le
vecteur concret.

Pour introduire une erreur contrôlée, les deux mémoires sont rationalisées sous la forme \(z_0=MK/\sqrt8\), et l'on ajoute \(1/10\) à la première coordonnée. Chaque bloc de douze composantes est ensuite transporté par les deux canaux bilatéraux. L'inverse par somme récupère les données de mémoire perturbées, et le décodeur de charge reçoit \(q=6263\) avec ces données, sans recevoir le registre recherché. La norme de l'erreur dans les mémoires originales est \(\sqrt8/10<r_q\). La normalisation isométrique conserve cette norme ; les canaux non normalisés facilitent seulement le calcul rationnel exact.

Dans cette instance, les candidats de la reconstruction par plancher et
plafond qui satisfont la charge sont \(924\). Le minimum de distance
quadratique rationalisée est \(1/100\), et le décodeur récupère
\eqref{x_reciprocidad:lib04:k-ejemplo}. L’unicité découle du rayon strict
\eqref{x_reciprocidad:lib04:radios}, non de la seule taille de l’énumération.
La lecture arithmétique récupérée est
\begin{equation}
 \kappa=
 \frac{234543140729659824621058914794146601}
      {999999999999999999999999999999999999},
 \label{x_reciprocidad:lib04:kappa-ejemplo}
\end{equation}
et la sélection incidencielle redevient \(\{4,6,9,10\}\), avec
le même seuil \(729\). L’exemple introduit le bruit avant le
transport. Les preuves de contractivité de l’adjoint couvrent également
des erreurs ultérieures arbitraires dont la norme conjointe appartient au
rayon déclaré.

La capacité opératoire réunie consiste à transporter et à retrouver un registre
engendré, à maintenir le conditionnement de ses réponses et à restituer
ses lectures arithmétiques et incidencielles avant d’appliquer des sélections
discontinues. Ses prémisses effectives sont la production préalable de
\(K\), le repère cyclique inversible, le bilan bilatéral, la
séparation discrète des mémoires et le théorème d’archivage. Le
registre de douze composantes conserve la portée de ses lecteurs ;
l’histoire enrichie complète maintient ses autres registres et
dépendances.
